% !TeX program = xelatex
% UTF-8; standalone XeLaTeX/xepersian document.
% Source: Google Drive / Notebook / kawasaki (ascending filename order).
% pic01.jpg: 20261018_192653057.jpg; 1XkA0FOD7GLNMG3YGfMi3YcVEbpIjdsIP
% pic02.jpg: 20261018_192657094.jpg; 1MTUAdgESDzT4Q7wAfOiM8ohujirpvWGO
% pic03.jpg: 20261018_192703082.jpg; 12x6-kDxlWsYzClULWJ-arw1-6BIQNmE5
\documentclass[11pt,a4paper]{article}
\usepackage[margin=20mm]{geometry}
\usepackage{amsmath,amssymb}
\usepackage{graphicx}
\usepackage{xcolor}
\usepackage{enumitem}
\usepackage{xepersian}
\settextfont{Amiri}
\setlatintextfont{DejaVu Serif}
\setlength{\parindent}{0pt}
\setlength{\parskip}{3pt}
\setlength{\emergencystretch}{3em}
\setlist[itemize]{itemsep=2pt,topsep=3pt}
\newcommand{\en}[1]{\lr{#1}}
\newcommand{\hi}[1]{\colorbox{yellow!35}{\strut #1}}
\newcommand{\up}{\ensuremath{\uparrow}}
\newcommand{\pagehead}[1]{\section*{صفحهٔ دست‌نویس #1}}
\begin{document}
\begin{center}
{\Large\bfseries ۸۸ ـ کاوازاکی}\\[3pt]
\en{Kawasaki disease}
\end{center}

\pagehead{۱}
\begin{itemize}
\item بیماری کاوازاکی: اتیولوژی ناشناخته؛ واسکولیت.
\item التهاب \hi{عروق کوچک و متوسط}؛ یک بیماری مولتی‌سیستمی.
\item منجر به تشکیل \hi{آنوریسم}.
\item دومین واسکولیت شایع در کودکان است: ۹ در هر \en{100,000} کودک زیر ۵ سال.
\item \hi{\en{F < M}}؛ شایع‌ترین سن: \hi{زیر ۵ سال}؛ پیک بین \hi{۲ـ۳ سالگی}.
\item بروز بیماری بین فوریه و مارس بالا می‌رود: \hi{زمستان}.
\item شایع‌ترین علت بیماری‌های قلبی اکتسابی در کودکان در آمریکا: کاوازاکی.
\end{itemize}

\subsection*{تظاهرات بالینی}
تظاهرات بالینی کاوازاکی در ۳ فاز بررسی می‌شود.

\subsubsection*{۱) \en{Acute}}
\begin{itemize}
\item \hi{۱ـ۲ هفته} طول می‌کشد.
\item شروع ناگهانی \hi{تب}؛ به دنبال تب، نشانه‌های مشخص:
\item اریتم ملتحمه، تغییرات مخاط، ترک لب، زبان توت‌فرنگی، لنفادنوپاتی گردنی، تورم دست و پا.
\item \hi{لنفادنوپاتی گردنی} در ۷۰\% موارد دیده می‌شود و باید \hi{بالای \en{1.5 cm}} باشد.
\item راش در ۸۰\% موارد؛ ممکن است در هر جایی باشد؛ اینگوینال و منتشر.
\item تحریک‌پذیری: بسیار تحریک‌پذیر؛ از صفات کاوازاکی.
\item درد شکم، هیدروپس کیسهٔ صفرا.
\item پلئوسیتوز \en{CSF}.
\item پیوری استریل.
\item آرتریت: مفاصل متوسط و بزرگ.
\item کاردیت: \hi{تاکی‌کاردی}، نارسایی شدید ناگهانی؛ \en{SoB}.
\item آنوریسم‌های شریان کرونری: بیشتر در کودکان بسیار کوچک.
\end{itemize}

\clearpage
\begin{center}{\large\bfseries ۸۸ ـ کاوازاکی}\end{center}
\pagehead{۲}
\subsubsection*{۲) \en{Subacute}}
\begin{itemize}
\item تا حدود \hi{\en{4w}} باقی می‌ماند.
\item تب به‌صورت تدریجی قطع می‌شود.
\item \hi{پوسته‌ریزی انگشتان دست و پا}.
\item عدد پلاکت تا الان \en{NL} بوده؛ \hi{افزایش} تا یک میلیون و بالاتر.
\item آشکار شدن آنوریسم عروق کرونر.
\end{itemize}

\textbf{ریسک فاکتورهای آنوریسم کرونری:}
\begin{itemize}
\item تب طول‌کشیده و \en{ESR} بالا، طولانی.
\item سن زیر یک سال یا بالای ۹ سال.
\item پاسخ ضعیف به درمان.
\item پسر بودن.
\end{itemize}

\subsubsection*{۳) \en{Convalescent}}
\begin{itemize}
\item از زمانی که علائم بالینی از بین برود شروع می‌شود: \hi{تا زمانی که \en{ESR} نرمال گردد}.
\item حدوداً \hi{۶ـ۸ هفته} پس از شروع بیماری.
\item \en{Linear nail creases}: \hi{\en{Beau Lines}}.
\end{itemize}

\subsection*{آزمایشگاه}
\begin{itemize}
\item سایر علل تب رد شوند: کشت ادرار، خون، \en{CXR}.
\item در فاز حاد: \en{WBC\up}، \en{CRP\up}، \hi{\en{ESR\up} (معمولاً بالای ۸۰)}.
\item \hi{\en{Plt} نرمال، یا کمی کاهش‌یافته است}.
\item در کاوازاکی قلب با اکوکاردیوگرافی بررسی می‌شود: در \hi{هفتهٔ ۱ـ۲} و \hi{هفتهٔ ۴ـ۶}.
\end{itemize}

\subsection*{تشخیص \en{KD}}
\hi{\en{FWLS} به‌مدت ۵ روز} و وجود حداقل \en{$4/5$} معیارها:
\begin{itemize}
\item \hi{کنژنکتیویت دوطرفه}؛ چرکی نیست.
\item حداقل یک تغییر در غشاهای مخاطی (لب‌ها ترک‌خورده، زبان توت‌فرنگی).
\item حداقل یک تغییر در اندام‌ها (اریتم پریفرال، پوسته‌ریزی انگشتان).
\item \hi{لنفادنوپاتی سینگل، سرویکال، بالای \en{1.5} سانتی‌متر}.
\item عدم وجود تشخیص دیگر.
\end{itemize}

\clearpage
\begin{center}{\large\bfseries ۸۸ ـ کاوازاکی}\end{center}
\pagehead{۳}
\subsection*{درمان \en{KD}}
\begin{itemize}
\item \hi{\en{IVIG} رکن اصلی درمان است}:
\begin{latin}
\noindent 2 g/kg over 12 h --- single dose.
\end{latin}
\item آسپرین با دوز ضدالتهابی:
\begin{latin}
\noindent 80--90 mg/kg/day --- divided every 6 hs.
\end{latin}
در فاز حاد بیماری.
\item زمانی که \hi{تب قطع شد}، دوز آسپرین را به آنتی‌ترومبوتیک کاهش می‌دهیم:
\begin{latin}
\noindent 3--5 mg/kg/day --- single dose.
\end{latin}
\item برای ۴ـ۶ هفته آسپرین می‌دهیم تا زمانی که اکو بگوید درگیری قلبی وجود ندارد.
\item اگر \hi{کاردیت} داشت یا \hi{تب بعد از ۲ دوز \en{IVIG} قطع نشد}:
\begin{latin}
\noindent Corticostroids or Infliximab.
\end{latin}
\end{itemize}

\begin{itemize}
\item اکثر موارد \en{KD} بدون عارضه درمان می‌شوند.
\item برخی: آنوریسم، \en{MI}، \en{SCD}؛ مننژیت آسپتیک.
\item \en{IVIG} ریسک آنوریسم کرونری را از ۲۵\% به ۵\% رسانده است.
\item التهاب سیستمیک شدید، ریسک آنوریسم کرونری را بالا می‌برد.
\end{itemize}

\subsection*{تشخیص آتیپیکال کاوازاکی (ناکامل)}
\begin{itemize}
\item در \hi{شیرخواران} شایع‌تر است.
\item تب حداقل به‌مدت ۵ روز + ۲ـ۳ مورد از کرایتریاها؛ خصوصاً در حضور آنوریسم کرونری.
\item شیرخوار زیر ۶ ماه + حداقل ۷ روز تب.
\end{itemize}

\subsection*{تشخیص افتراقی‌های کاوازاکی}
\begin{itemize}
\item عفونت.
\item التهاب: \en{JIA}، بهجت، پلی‌آرتریت ندوزا.
\item ازدیاد حساسیت: \en{SJS}، واکنش دارویی.
\end{itemize}

\clearpage
\section*{تصاویر اصلی}
\IfFileExists{images/pic01.jpg}{\def\imagebase{images/}}{%
\def\imagebase{pediatrics/ped-rheumatology/kawasaki disease/images/}}
\begin{center}
\textbf{تصویر ۱}\par
\includegraphics[width=\linewidth,height=0.88\textheight,keepaspectratio]{\imagebase pic01.jpg}
\end{center}
\clearpage
\begin{center}
\textbf{تصویر ۲}\par
\includegraphics[width=\linewidth,height=0.88\textheight,keepaspectratio]{\imagebase pic02.jpg}
\end{center}
\clearpage
\begin{center}
\textbf{تصویر ۳}\par
\includegraphics[width=\linewidth,height=0.88\textheight,keepaspectratio]{\imagebase pic03.jpg}
\end{center}
\end{document}
